The system uses 8 DynamoDB tables. Before creating them, there is one
fundamental difference from relational databases worth understanding.

\vspace{0.5em}\noindent\textbf{Design around access
patterns}\par

With a relational database you normalise the data first and write
queries later. With DynamoDB it is the other way round: you decide how
you will query first, then design tables to serve exactly those queries.

Here, the first six tables only ever need a lookup by a single key. The
last two are different:

\begin{itemize}
\tightlist
\item
  \textbf{Reviews}: needs \emph{all reviews for one product}
\item
  \textbf{Orders}: needs \emph{all orders for one user}
\end{itemize}

With a single key, fetching those records would require a \textbf{Scan}
that reads the entire table and filters. That is slow and expensive,
since DynamoDB bills by the volume of data read.

The solution is a \textbf{composite key} made of a partition key and a
sort key. That lets you use \textbf{Query}, which reads only the
partition you need.

{\def\LTcaptype{none} % do not increment counter
\begingroup
\small
\setlength{\tabcolsep}{3pt}
\renewcommand{\arraystretch}{1.15}
\begin{longtable}[]{@{}llll@{}}
\toprule\noalign{}
Table & Partition key & Sort key & Contents \\
\midrule\noalign{}
\endhead
\bottomrule\noalign{}
\endlastfoot
Products & \texttt{id} & --- & Product information \\
Categories & \texttt{id} & --- & Categories \\
Users & \texttt{id} & --- & Accounts, hashed passwords \\
Sessions & \texttt{token} & --- & Login sessions \\
Carts & \texttt{userId} & --- & Shopping carts \\
Vouchers & \texttt{code} & --- & Discount codes \\
Reviews & \texttt{productId} & \texttt{id} & Product reviews \\
Orders & \texttt{userId} & \texttt{id} & Orders \\
\end{longtable}
\endgroup
}

\textbf{Production Note:} DynamoDB is highly suitable for high-speed
transactional workflows (like shopping carts and orders). However, for
complex product search and filtering queries (such as searching by price
range, color, or free-text keywords), you should consider syncing data
to \textbf{Amazon OpenSearch} to process them rather than relying solely
on DynamoDB.

\vspace{0.5em}\noindent\textbf{Step 1. Create a table in the
console}\par

Do one by hand first to get familiar with the interface:

\begin{enumerate}
\def\labelenumi{\arabic{enumi}.}
\tightlist
\item
  Open the
  \href{https://ap-southeast-1.console.aws.amazon.com/dynamodbv2/home?region=ap-southeast-1}{DynamoDB
  console}
\item
  Choose \textbf{Create table}
\item
  Fill in:

  \begin{itemize}
  \tightlist
  \item
    \textbf{Table name}: \texttt{Products}
  \item
    \textbf{Partition key}: \texttt{id}, type \textbf{String}
  \item
    Leave Sort key empty
  \item
    \textbf{Table settings}: choose \textbf{Customize settings}
  \item
    \textbf{Table class}: DynamoDB Standard
  \item
    \textbf{Capacity mode}: \textbf{On-demand}
  \end{itemize}
\item
  Choose \textbf{Create table}
\end{enumerate}

\begin{figure}
\centering
\pandocbounded{\includegraphics[keepaspectratio,alt={Create table}]{5-Workshop/5.4/create-table.png}}
\caption{Create table}
\end{figure}

\textbf{On-demand} means you pay for actual reads and writes rather than
pre-provisioned capacity. For a project with unknown traffic this is the
safe choice: you neither overpay nor get throttled.

\vspace{0.5em}\noindent\textbf{Step 2. Create the remaining seven tables via
CLI}\par

Creating seven tables by hand is tedious. Use this instead:

\begin{Shaded}
\begin{Highlighting}[]
\ControlFlowTok{for}\NormalTok{ T }\KeywordTok{in}\NormalTok{ Categories Users Sessions Carts Vouchers}\KeywordTok{;} \ControlFlowTok{do}
  \ControlFlowTok{case} \VariableTok{$T} \KeywordTok{in}
    \SpecialStringTok{Sessions}\KeywordTok{)} \VariableTok{K}\OperatorTok{=}\NormalTok{token }\ControlFlowTok{;;}
    \SpecialStringTok{Carts}\KeywordTok{)}    \VariableTok{K}\OperatorTok{=}\NormalTok{userId }\ControlFlowTok{;;}
    \SpecialStringTok{Vouchers}\KeywordTok{)} \VariableTok{K}\OperatorTok{=}\NormalTok{code }\ControlFlowTok{;;}
    \PreprocessorTok{*}\KeywordTok{)}        \VariableTok{K}\OperatorTok{=}\NormalTok{id }\ControlFlowTok{;;}
  \ControlFlowTok{esac}
  \ExtensionTok{aws}\NormalTok{ dynamodb create{-}table }\AttributeTok{{-}{-}table{-}name} \VariableTok{$T} \DataTypeTok{\textbackslash{}}
    \AttributeTok{{-}{-}attribute{-}definitions}\NormalTok{ AttributeName=}\VariableTok{$K}\NormalTok{,AttributeType=S }\DataTypeTok{\textbackslash{}}
    \AttributeTok{{-}{-}key{-}schema}\NormalTok{ AttributeName=}\VariableTok{$K}\NormalTok{,KeyType=HASH }\DataTypeTok{\textbackslash{}}
    \AttributeTok{{-}{-}billing{-}mode}\NormalTok{ PAY\_PER\_REQUEST}
\ControlFlowTok{done}
\end{Highlighting}
\end{Shaded}

The two composite-key tables are created separately:

\begin{Shaded}
\begin{Highlighting}[]
\ExtensionTok{aws}\NormalTok{ dynamodb create{-}table }\AttributeTok{{-}{-}table{-}name}\NormalTok{ Reviews }\DataTypeTok{\textbackslash{}}
  \AttributeTok{{-}{-}attribute{-}definitions}\NormalTok{ AttributeName=productId,AttributeType=S AttributeName=id,AttributeType=S }\DataTypeTok{\textbackslash{}}
  \AttributeTok{{-}{-}key{-}schema}\NormalTok{ AttributeName=productId,KeyType=HASH AttributeName=id,KeyType=RANGE }\DataTypeTok{\textbackslash{}}
  \AttributeTok{{-}{-}billing{-}mode}\NormalTok{ PAY\_PER\_REQUEST}

\ExtensionTok{aws}\NormalTok{ dynamodb create{-}table }\AttributeTok{{-}{-}table{-}name}\NormalTok{ Orders }\DataTypeTok{\textbackslash{}}
  \AttributeTok{{-}{-}attribute{-}definitions}\NormalTok{ AttributeName=userId,AttributeType=S AttributeName=id,AttributeType=S }\DataTypeTok{\textbackslash{}}
  \AttributeTok{{-}{-}key{-}schema}\NormalTok{ AttributeName=userId,KeyType=HASH AttributeName=id,KeyType=RANGE }\DataTypeTok{\textbackslash{}}
  \AttributeTok{{-}{-}billing{-}mode}\NormalTok{ PAY\_PER\_REQUEST}
\end{Highlighting}
\end{Shaded}

\vspace{0.5em}\noindent\textbf{Step 3. Verify}\par

\begin{Shaded}
\begin{Highlighting}[]
\ExtensionTok{aws}\NormalTok{ dynamodb list{-}tables}
\end{Highlighting}
\end{Shaded}

The output should list all 8 tables.

\vspace{0.5em}\noindent\textbf{Step 4. Seed sample data}\par

Load sample categories and products so the system has something to
display:

\begin{Shaded}
\begin{Highlighting}[]
\BuiltInTok{cd}\NormalTok{ backend}
\ExtensionTok{node}\NormalTok{ scripts/seed{-}data.mjs}
\end{Highlighting}
\end{Shaded}

Check the record count in the Products table:

\begin{Shaded}
\begin{Highlighting}[]
\ExtensionTok{aws}\NormalTok{ dynamodb scan }\AttributeTok{{-}{-}table{-}name}\NormalTok{ Products }\AttributeTok{{-}{-}select}\NormalTok{ COUNT}
\end{Highlighting}
\end{Shaded}
